\documentclass[conference]{IEEEtran}

\usepackage{cite}
\usepackage{amsmath,amsfonts}
\usepackage{graphicx}
\usepackage{url}

\begin{document}

\title{CertiLock Campus: A Hybrid Blockchain and Cryptographic Hashing Framework for Secure Academic and Skill Certificate Verification}

\author{
\IEEEauthorblockN{Bhushan Lilhare}
\IEEEauthorblockA{
Department of Computer Science and Engineering\\
GHRCE, Nagpur, India \\
Email: your-email@example.com}
}

\maketitle

\begin{abstract}
The rapid increase in digital academic records has led to significant challenges in ensuring the authenticity and integrity of certificates. Traditional verification methods are often manual, time-consuming, and vulnerable to forgery. This paper presents CertiLock Campus, a hybrid framework that leverages blockchain technology and cryptographic hashing to provide a secure, tamper-proof system for academic and skill certificate verification. The proposed system generates a unique cryptographic hash (SHA-256) for each certificate and stores it on a blockchain network, ensuring immutability and transparency. Additionally, QR codes are integrated to enable real-time verification by third parties. The system utilizes a hybrid storage architecture combining blockchain for verification and cloud-based storage for efficient data management. Experimental results demonstrate that the proposed system significantly reduces verification time while enhancing security and trust. This framework can be effectively adopted by educational institutions to combat certificate fraud and streamline verification processes.
\end{abstract}

\begin{IEEEkeywords}
Blockchain, Academic Certificates, Cryptographic Hashing, QR Code, Smart Contracts, Verification
\end{IEEEkeywords}

\section{Introduction}

With the digital transformation of educational systems, academic certificates and skill credentials are increasingly issued and stored electronically. However, this shift has also led to a rise in fraudulent activities such as certificate forgery and unauthorized modifications. Traditional verification processes rely heavily on manual checks or centralized databases, which are inefficient, time-consuming, and prone to security vulnerabilities.

Blockchain technology has emerged as a promising solution due to its decentralized, immutable, and transparent nature. By storing verification data on a distributed ledger, blockchain ensures that once a certificate is recorded, it cannot be altered or tampered with.

This paper proposes CertiLock Campus, a hybrid blockchain-based framework integrated with cryptographic hashing and QR code technology to provide a secure and efficient solution for academic and skill certificate verification. The system not only ensures data integrity but also enables instant verification by employers, institutions, and other stakeholders.

\section{Literature Review}
Several studies have explored the application of blockchain in academic credential verification. Existing systems primarily focus on storing certificates directly on the blockchain or using decentralized identifiers.

Some research proposes blockchain-based certificate systems that ensure immutability but suffer from scalability and storage limitations. Other approaches integrate IPFS (InterPlanetary File System) for decentralized storage, improving efficiency but introducing complexity in system design.

Additionally, QR code-based verification systems have been developed to provide quick access to certificate details; however, these systems often rely on centralized databases, making them vulnerable to tampering.

Despite these advancements, most existing solutions lack a hybrid architecture that combines blockchain security with efficient storage and real-time verification mechanisms. The proposed CertiLock Campus system addresses these limitations by integrating blockchain, cryptographic hashing, and cloud storage into a unified framework.

\section{Proposed System}

\subsection{System Architecture}

The system consists of the following modules:
\begin{itemize}
\item Admin Module: Issues certificates
\item Student Module: Accesses certificates
\item Verifier Module: Verifies certificates using QR codes
\item Blockchain Layer: Stores certificate hashes
\item Cloud Storage: Stores certificate data
\end{itemize}

\subsection{Working Process}

\begin{enumerate}
\item Certificate is generated by admin
\item SHA-256 hash is created
\item Hash is stored on blockchain
\item Certificate stored in cloud
\item QR code generated
\item Verifier scans QR code
\item Hash is compared
\item Validity is confirmed
\end{enumerate}

\section{Methodology and Implementation}

\subsection{Blockchain Layer}
Ethereum/Polygon blockchain is used to store certificate hashes via smart contracts.

\subsection{Cryptographic Layer}
SHA-256 hashing ensures integrity and generates unique fingerprints.

\subsection{Storage Layer}
Firebase/IPFS is used to store certificate files efficiently.

\subsection{Application Layer}
Frontend developed using React/Android, backend using Node.js and Firebase.

\section{Results and Discussion}

The system demonstrates:
\begin{itemize}
\item Fast verification (few seconds)
\item High security due to immutability
\item Reduced fraud
\item Improved scalability using hybrid model
\end{itemize}

Compared to traditional systems, the proposed solution significantly improves efficiency and trust.

\section{Advantages}

\begin{itemize}
\item Tamper-proof certificates
\item Real-time verification
\item Decentralized trust
\item Reduced manual effort
\item Scalable architecture
\end{itemize}

\section{Conclusion}

This paper presented CertiLock Campus, a hybrid blockchain-based framework for secure academic and skill certificate verification. By integrating blockchain, cryptographic hashing, and QR code technology, the system ensures tamper-proof, transparent, and efficient verification. The hybrid architecture enhances scalability while maintaining high security standards. The proposed system has the potential to transform the way educational institutions manage and verify credentials.

\section{Future Scope}

\begin{itemize}
\item Integration with national education systems
\item AI-based fraud detection
\item Self-sovereign identity (SSI)
\item Mobile app enhancements
\end{itemize}

\begin{thebibliography}{00}

\bibitem{b1} S. Nakamoto, ``Bitcoin: A Peer-to-Peer Electronic Cash System,'' 2008.

\bibitem{b2} M. Sharples and J. Domingue, ``The Blockchain and Kudos: A Distributed System for Educational Record,'' 2016.

\bibitem{b3} Z. Zheng et al., ``Blockchain Challenges and Opportunities,'' IEEE, 2017.

\bibitem{b4} K. Christidis and M. Devetsikiotis, ``Blockchains and Smart Contracts for the Internet of Things,'' IEEE, 2016.

\bibitem{b5} A. Grech and A. Camilleri, ``Blockchain in Education,'' European Commission, 2017.

\end{thebibliography}

\end{document}